\section{Instruccions d’instal·lació i execució}
\label{annex:installacio}

Aquest annex descriu els passos necessaris per recuperar el codi font, instal·lar les dependències i executar els diferents components del prototip. L’objectiu és facilitar la reproducció de l’entorn utilitzat durant el desenvolupament i l’avaluació del sistema.

\subsection{Requisits previs}

Abans d’executar el projecte, és necessari disposar dels components següents:

\begin{itemize}
    \item Node.js i npm;
    \item MetaMask Flask;
    \item Ollama;
    \item el model local \texttt{llama3.2:latest};
    \item Git;
    \item un navegador compatible amb MetaMask Flask;
    \item accés a un endpoint RPC de Sepolia;
    \item ETH de prova de la xarxa Sepolia.
\end{itemize}

Les versions exactes utilitzades durant l’avaluació es recullen a la Taula~\ref{tab:configuracio_avaluacio}.

\subsection{Descàrrega del repositori}

El codi font del projecte es pot obtenir des del repositori de GitHub mitjançant l’ordre següent:

\begin{lstlisting}[language=bash,breaklines=true,breakatwhitespace=false]
git clone https://github.com/marccarot0202/TFG_MarcCarotFigueres.git
cd TFG_MarcCarotFigueres
\end{lstlisting}

La versió utilitzada en l'avaluació correspon a l'estat del repositori dels dies 17 i 18 de juliol de 2026. Abans de publicar la versió final es recomana crear una etiqueta Git associada a l'entrega.

\subsection{Instal·lació de les dependències}

Les dependències generals del projecte s’instal·len des de l’arrel del repositori:

\begin{lstlisting}[language=bash]
yarn install
\end{lstlisting}

El backend i els contractes disposen també dels seus propis fitxers de dependències:

\begin{lstlisting}[language=bash]
cd backend
npm install

cd ../smart-contracts
npm install
\end{lstlisting}

\subsection{Instal·lació del model local}

El model utilitzat pel prototip es descarrega mitjançant Ollama:

\begin{lstlisting}[language=bash]
ollama pull llama3.2
\end{lstlisting}

Es pot comprovar que el model està disponible amb:

\begin{lstlisting}[language=bash]
ollama list
\end{lstlisting}

El servei local d’Ollama ha d’estar accessible a:

\begin{center}
\texttt{http://localhost:11434}
\end{center}

\subsection{Configuració de les variables d’entorn}

Les credencials i les dades sensibles no s’emmagatzemen directament al repositori. Per executar el desplegament o les consultes sobre Sepolia és necessari crear un fitxer local de variables d’entorn.

Un exemple simplificat de configuració és el següent:

\begin{lstlisting}
SEPOLIA_RPC_URL=URL_DEL_PROVEIDOR_RPC
PRIVATE_KEY=CLAU_PRIVADA_DEL_COMPTE_DE_PROVES
ETHERSCAN_API_KEY=CLAU_DE_L_API
\end{lstlisting}

Aquest fitxer no s’ha d’incorporar al control de versions ni publicar al repositori.

\subsection{Execució dels components}

El prototip està format per diversos serveis que s’han d’executar simultàniament. En l'entorn Windows utilitzat per a l'avaluació, la forma recomanada és executar \texttt{Iniciador\_Portatil.bat} des de l'arrel. Aquest script selecciona Node.js 20.11.1 i obre Ollama, el backend i el \textit{workspace} de la DApp i el Snap en l'ordre correcte.

\begin{lstlisting}[language=bash]
Iniciador_Portatil.bat
\end{lstlisting}

Les ordres següents descriuen els components per separat amb finalitat de reproducció i diagnòstic; no s'han d'executar simultàniament amb l'iniciador, perquè es produirien conflictes de ports.

\subsubsection{Backend local}

El backend d’anàlisi s’executa al port \texttt{3000}:

\begin{lstlisting}[language=bash]
cd backend
node server.js
\end{lstlisting}

Es pot comprovar el seu estat accedint als endpoints:

\begin{center}
\texttt{http://localhost:3000/health}
\end{center}

\begin{center}
\texttt{http://localhost:3000/stats}
\end{center}

\subsubsection{Snap}

El servidor de desenvolupament del Snap s’executa al port \texttt{8080}:

\begin{lstlisting}[language=bash]
yarn workspace snap start
\end{lstlisting}

\subsubsection{DApp local}

La DApp de proves s’executa al port \texttt{8000}. L’ordre exacta depèn de l’estructura del repositori:

\begin{lstlisting}[language=bash]
yarn workspace site start
\end{lstlisting}

Un cop iniciada, la interfície queda disponible a:

\begin{center}
\texttt{http://localhost:8000}
\end{center}

\subsection{Ordre recomanat d’inicialització}

Per preparar correctament l’entorn, es recomana seguir aquest ordre:

\begin{enumerate}
    \item iniciar Ollama;
    \item iniciar el backend local;
    \item iniciar el servidor del Snap;
    \item iniciar la DApp;
    \item obrir MetaMask Flask;
    \item instal·lar o reconnectar el Snap;
    \item seleccionar la xarxa Sepolia;
    \item comprovar l’estat dels serveis des de la DApp;
    \item iniciar una operació de prova.
\end{enumerate}
